% Generated draft: requires mathematical and source review.
\documentclass[11pt]{article}
\usepackage[margin=25mm]{geometry}
\usepackage{fontspec}
\setmainfont{Latin Modern Roman}
\usepackage{amsmath,amssymb,mathtools,unicode-math}
\setmathfont{Latin Modern Math}
\usepackage{enumitem,needspace,xurl,hyperref}
\setlength{\parindent}{0pt}
\setlength{\parskip}{5pt}
\setlength{\emergencystretch}{3em}
\newcommand{\N}{\mathbb{N}}
\newcommand{\Z}{\mathbb{Z}}
\newcommand{\Q}{\mathbb{Q}}
\newcommand{\R}{\mathbb{R}}
\newcommand{\C}{\mathbb{C}}
\newcommand{\F}{\mathbb{F}}
\newcommand{\A}{\mathbb{A}}
\newcommand{\PP}{\mathbb{P}}
\newcommand{\E}{\mathbb{E}}
\newcommand{\eps}{\varepsilon}
\newcommand{\dd}{\,\mathrm{d}}
\begin{document}
\section*{UnsolvedMath ID: 800507}
\textbf{Source problem code:} catalogue entry
\subsection{Definitions and mathematical statement}
In the conventional real symmetric interpretation, a symmetric positive definite linear system is
\[
Ax=b,\qquad A\in\R^{n\times n},\quad b\in\R^n,\quad n\geq 1,
\]
where \(A^{\mathsf T}=A\) and \(v^{\mathsf T}Av>0\) for every nonzero \(v\in\R^n\). The source does not explicitly specify the scalar domain; this real formulation is an interpretation requiring confirmation. CG denotes the conjugate gradient algorithm. Here the Lanczos algorithm is to be used as a method for solving the linear system, but the source does not specify how its iterates are constructed.

\textbf{Source question.} Determine whether CG and Lanczos behave similarly in finite precision arithmetic when solving a symmetric positive definite linear system, using implementations satisfying the source's requirement of ``standard good implementations'' of both methods. Preserve the source's qualification that the answer depends on the implementations.

\textbf{Unresolved formulation.} The source does not define ``behave similarly'' or specify the implementations, the arithmetic model, initialization, stopping criteria, or quantifiers over systems and precision. Consequently, it supplies a research question rather than a fully specified mathematical proposition. A precise comparison requires these choices and a mathematical criterion for similarity; no particular criterion or universal equivalence assertion is imposed here.
\subsection{Short English statement}
The question asks whether conjugate gradient and Lanczos behave similarly when solving symmetric positive definite linear systems in finite precision. The answer is explicitly implementation dependent. The supplied record leaves both the comparison criterion and the meaning of ``standard good implementations'' unspecified, so a precise theorem or counterexample target requires further formulation.
\subsection{Sources}
\begin{itemize}
\item Local UnsolvedMath registry: permanent ID \(800507\), exact source problem code \texttt{catalogue entry}, title \emph{Linear Systems and Eigenvalue Problems --- Problem 16}. The registry attributes the record to Amsel, Baumann et al., with proposed year 2026, and labels it \texttt{open}. This is a supplied status label, not an independently verified mathematical assessment.
\item Supplied source: Amsel, Baumann et al., \emph{Linear Systems and Eigenvalue Problems: Open Questions from a Simons Workshop} (2026), Problem 16, \url{https://arxiv.org/abs/2602.05394}. The supplied background records access on 2026-10-16 and extraction by \texttt{arxiv\_latex}; these are supplied provenance statements, not independently verified access or source verification.
\end{itemize}
\end{document}
